\documentclass[11pt]{article}
\usepackage[margin=25mm]{geometry}
\usepackage{fontspec}
\setmainfont{Latin Modern Roman}
\usepackage{amsmath,amssymb,mathtools}
\usepackage{unicode-math}
\setmathfont{Latin Modern Math}
\usepackage{xurl,hyperref}
\hypersetup{colorlinks=true,urlcolor=blue}
\setcounter{secnumdepth}{0}
\setlength{\parindent}{0pt}
\setlength{\parskip}{5pt}
\setlength{\emergencystretch}{3em}
\newcommand{\sref}[2]{\hyperlink{source:#1:#2}{[S#2]}}
\newcommand{\cataloguescope}[1]{\paragraph{Recorded target.}#1}
\newcommand{\R}{\mathbb R}
\newcommand{\C}{\mathbb C}
\newcommand{\N}{\mathbb N}
\newcommand{\Z}{\mathbb Z}
\newcommand{\Q}{\mathbb Q}
\newcommand{\F}{\mathbb F}
\newcommand{\eps}{\varepsilon}
\begin{document}
\section{Open Structures in Compactified Imaginary Liouville Theory}\label{statement}

\subsection{Definitions and mathematical statement}

\paragraph{Shared notation.}$\mathbb N=\{1,2,\ldots\}$; the standard number fields have their usual meanings. All additional conventions are stated locally.


Imaginary Liouville theory uses an imaginary exponential coupling of a scalar field, with parameter $\beta$. Compactification identifies field values periodically. A logarithmic conformal theory may have nonsemisimple dilatation action and logarithmic correlations. CILT refers to the compactified imaginary Liouville theory defined in the report. BPZ minimal models are conformal field theories with central charge $c=1-6(p-q)^2/(pq)$ for coprime integers $p>q$. Felder's BRST reduction extracts the minimal-model state spaces from Coulomb-gas representations.

\paragraph{Questions recovered from the primary report.}The report's five questions ask: (i) how BPZ minimal models are related to CILT; (ii) whether Felder's BRST reduction has a path-integral formulation; (iii) whether the non-self-adjoint CILT Hamiltonian and its Jordan cells yield a logarithmic conformal field theory; (iv) how to construct CILT for $\beta\in[\sqrt2,2)$, where Gaussian multiplicative chaos renormalization is insufficient; and (v) whether there is a probabilistic theory when $\beta^2\notin\mathbb Q$.

\paragraph{Catalogue wording retained for traceability.}
Are compactified imaginary Liouville models related to CILT, do they admit a path-integral formulation, are they logarithmic CFTs, and is there a probabilistic theory when $\beta^2\notin\mathbb{Q}$?

\paragraph{Scope and source qualification.}The catalogue condenses five primary-source questions to four and omits the $\beta\in[\sqrt2,2)$ construction question. Its first relationship is with BPZ minimal models, and its path-integral question concerns Felder's BRST reduction. The report already constructs CILT in a smaller rational-coupling regime. The recovered questions above correct these ambiguities; the recorded catalogue wording below is retained for traceability. No positivity of the complex formal path-integral weight is assumed.
\subsection{Short English statement}
Relate compactified imaginary Liouville theory to minimal models, formulate Felder's reduction by a path integral, characterize its logarithmic structure and construct the remaining coupling regimes.
\subsection{Research attempt}
\paragraph{Model, date and investigated scope.}
OpenAI GPT-6 Astra (configured model \texttt{gpt-6-astra}), ultra reasoning; generated 11 October 2026. Derived the local imaginary-Gaussian second-moment threshold and the logarithmic effect of an assumed Jordan block; all five recovered CILT questions remain.

\paragraph{Sources and literature check.}
Read the full supplied definition, catalogue wording, source qualifications and listed references for this question. No fresh status or novelty search was performed. The argument below proves its auxiliary statements directly and does not import an unverified external theorem.

\paragraph{The renormalization threshold in a local test.} Assume a regularized real Gaussian field on a bounded planar region has covariance \(C_\epsilon(x,y)\) approaching \(-\log|x-y|+O(1)\) off the diagonal. For \(Z_\epsilon(x)=\exp(i\beta X_\epsilon(x)+\beta^2C_\epsilon(x,x)/2)\), direct Gaussian integration gives
\[
 \mathbb E\bigl[Z_\epsilon(x)\overline{Z_\epsilon(y)}\bigr]
 =\exp(\beta^2 C_\epsilon(x,y)).
\]
The candidate second-moment kernel is therefore comparable locally to \(|x-y|^{-\beta^2}\). Its radial integral near the diagonal is \(2\pi\int_0^1r^{1-\beta^2}dr\), finite exactly when \(\beta^2<2\), logarithmically divergent at equality. This diagnoses why a plain second-moment construction stops at \(\sqrt2\); it does not rule out other renormalizations beyond that point. Rationality of \(\beta^2\) plays no role in this local integrability test.
\paragraph{Jordan cells and logarithms.} If a dilatation operator has a block \(H=hI+N\) with \(N^2=0\) but \(N\ne0\), then for \(r>0\)
\[
 r^{-H}=r^{-h}\bigl(I-(\log r)N\bigr).
\]
This is the exact truncated exponential series. A correlation sees a logarithm only if its chosen bra and ket couple to \(N\). Non-self-adjointness alone proves no such block: \(\operatorname{diag}(i,2i)\) is non-self-adjoint and diagonalizable.
\paragraph{All five source questions remain.} These diagnostics neither identify BPZ minimal-model state spaces inside CILT nor formulate Felder's BRST reduction as a path integral. They do not prove actual CILT Jordan cells/logarithmic correlation axioms, construct \(\beta\in[\sqrt2,2)\), or construct the irrational-coupling probabilistic theory. The first analytic gap is a controlled non-Gaussian/renormalized CILT limit beyond the displayed cutoff test, together with the relevant state-space identifications. No positive interpretation of the complex path-integral weight is assumed.

\paragraph{Verification and adversarial audit.}
Checked complex conjugation in the covariance formula, the beta=square-root-two endpoint, rationality independence and the requirement of a nonzero Jordan matrix element. Model self-audit only; no independent human mathematical verification. Applied contribution prompt A before drafting and B afterwards. Checked complex conjugation in the covariance formula, the beta=square-root-two endpoint, rationality independence and the requirement of a nonzero Jordan matrix element. First remaining gap: Construct the required renormalized CILT limits and identify their BRST/state-space structure beyond these local algebraic tests. The inherited source status is preserved; original-target completion remains zero.

Coordinator checked the conjugate second-moment kernel, planar threshold and nilpotent-block logarithm. This additional check was model-assisted, not independent human verification.

\paragraph{Tools, process and contributions.}
Python 3 and local file tools used to read definitions and serialize manuscripts. Calculations are symbolic reasoning; no numerical experiment or theorem prover. Contribution prompt A was followed by prompt B. GPT-6 Astra with ultra reasoning drafted the research and performed its stated self-audit; any additional coordinator contribution is identified in the verification and audit. A GPT-6-based Codex coordinator assembled the manuscript and checked repository integration; its exact variant and reasoning configuration were not exposed. No human mathematical contribution, independent human proof check or formal proof certification is supplied. Complete derivations and the consulted source references are included here; private reasoning transcripts are not published.

\paragraph{Final status and first remaining gap.}
\textbf{UNRESOLVED}. Construct the required renormalized CILT limits and identify their BRST/state-space structure beyond these local algebraic tests.

\section{Completion Estimate}
\noindent\textbf{COMPLETION: 0\%}
This is an elementary auxiliary calculation or an explicit test of a restricted strategy. No new portion of the full source target is established.

\subsection{Sources}
\begin{itemize}
\item[\hypertarget{source:statement:1}{S1}] ULAM.AI and the UnsolvedMath Contributors, UnsolvedMath, catalogue entry. Catalogue curation under CC BY 4.0; definitions and English exposition added locally.
\url{https://huggingface.co/datasets/ulamai/UnsolvedMath}.
\item[\hypertarget{source:statement:2}{S2}] Statistical Physics and Random Surfaces (2024), Oberwolfach Reports, DOI 10.4171/owr/2024/25, pp. 1387--1444. Original report and mathematical context. Incorporated report material remains under the source CC BY-SA 4.0 terms unless otherwise noted.
\url{https://doi.org/10.4171/owr/2024/25}.
\end{itemize}
\label{end:statement}
\end{document}
